\documentclass[12pt,letterpaper]{article}



%%
%% Required packages:
%%
\usepackage{etex}
\usepackage{amsmath}
\usepackage{amssymb}
\usepackage{amstext}
\usepackage{amsthm}
\usepackage{fancyhdr,fancybox}
\usepackage{mathrsfs} % need for \mathscr command
\usepackage{multicol}
\usepackage[normalem]{ulem}
%\usepackage{pictex,pictexplus}
% \usepackage{pictexwd}
\usepackage{rotating}
\usepackage{url}
\usepackage{xfrac}
\usepackage{booktabs}
\usepackage{color,soul}
\usepackage{comment}

\usepackage{hyperref}
\hypersetup{
    colorlinks=true,     % false: boxed links; true: colored links
    linkcolor=blue,      % color of internal links
    citecolor=blue,      % color of links to bibliography
    filecolor=blue,      % color of file links
    urlcolor=blue        % color of external links
}


\usepackage{tikz}
\usetikzlibrary{backgrounds,calc}

%%
%% Common formatting commands
%%
\setlength{\parindent}{0in}
\setlength{\textwidth}{7in}
\setlength{\evensidemargin}{-0.25in}
\setlength{\oddsidemargin}{-0.25in}
\setlength{\parskip}{.5\baselineskip}
\setlength{\topmargin}{-0.5in}
\setlength{\textheight}{9in}

%               Problem and Part
\newcounter{problemnumber}
\newcounter{partnumber}
\newcounter{subpartnumber}
\newcommand{\Problem}{%
  \refstepcounter{problemnumber}%
  \setcounter{partnumber}{0}%
  \setcounter{subpartnumber}{0}%
  \item[\makebox{\hfil\textbf{\theproblemnumber.}\hfil}]%
  }
\newcommand{\InProblem}[2][1.6]{%
  \refstepcounter{problemnumber}%
  \setcounter{partnumber}{0}%
  \setcounter{subpartnumber}{0}%
  \textbf{\theproblemnumber.}\ \ \parbox[t]{#1in}{#2}%
}
\newcommand{\InSmallProblem}[1]{\InProblem[1.05]{#1}}
\newcommand{\NoProblem}[1]{%
  \mbox{\hfil\textbf{#1}\hfil}%
  }
\newcommand{\UnProblem}{%
  \refstepcounter{problemnumber}%
  \setcounter{partnumber}{0}%
  \setcounter{subpartnumber}{0}%
  \mbox{\hfil\textbf{\theproblemnumber.}\hfil}%
  }
\newcommand{\InFlexProblem}[1]{%
  \refstepcounter{problemnumber}%
  \setcounter{partnumber}{0}%
  \setcounter{subpartnumber}{0}%
  \textbf{\theproblemnumber.}\ \ #1%
}
%%  Part:
\newcommand{\Part}{%
  \renewcommand{\thepartnumber}{(\alph{partnumber})}%
  \refstepcounter{partnumber}
  \setcounter{subpartnumber}{0}%
  \item[(\alph{partnumber})]%
}
\newcommand{\InPart}[2][1.75]{%
  \renewcommand{\thepartnumber}{(\alph{partnumber})}%
  \refstepcounter{partnumber}%
  \setcounter{subpartnumber}{0}%
  (\alph{partnumber})\ \ \parbox[t]{#1in}{#2}%
}
\newcommand{\UnPart}{%
  \refstepcounter{partnumber}%
  \renewcommand{\thepartnumber}{(\alph{partnumber})}%
  \setcounter{subpartnumber}{0}%
  (\alph{partnumber})%
}
\newcommand{\InLargePart}[1]{\InPart[3.05]{#1}}
\newcommand{\InFullPart}[1]{\InPart[6.2]{#1}}
\newcommand{\InSmallPart}[1]{\InPart[1.05]{#1}}
%% SubPart:
\newcommand{\SubPart}{%
  \refstepcounter{subpartnumber}%
  \item[(\roman{subpartnumber})]%
  }
\newcommand{\InSubPart}[2][2.25]{%
  \stepcounter{subpartnumber}%
  (\roman{subpartnumber})\ \ \parbox[t]{#1in}{#2}%
}
\newcommand{\InSmallSubPart}[1]{\InSubPart[1.5]{#1}}
\newcommand{\InTinySubPart}[1]{%
  \stepcounter{subpartnumber}%
  (\roman{subpartnumber})\ \ \makebox{#1}%
}
\newcommand{\InMySubPart}[2][2.25]{%
  \stepcounter{subpartnumber}%
  \parbox[t]{#1in}{\makebox[0.3in][l]{(\roman{subpartnumber})}\ \ {#2}}%
}
\newcommand{\TrueFalse}{\textbf{T}\ \ \textbf{F}\ \ }
\newcommand{\TFPart}[1]{% 
  \stepcounter{partnumber}%
  \setcounter{subpartnumber}{0}%
  \item[(\alph{partnumber})] \TrueFalse\parbox[t]{5.5in}{#1}}

\newcommand{\Points}[1]{(#1 points) \ }
\newcommand{\Point}[1]{(#1 point) \ }


%% For Answers:
\newcounter{answernumber}
\newcommand{\Answer}{%
  \refstepcounter{answernumber}%
  \setcounter{partnumber}{0}%
  \setcounter{subpartnumber}{0}%
  \item[\makebox{\hfil\textbf{\theanswernumber.}\hfil}]%
  }


% Setting up the headers for the first page
%\chead{} % will default to what is typed in the file.
\fancyhead[L]{\textsc{Math 22a}}
\fancyhead[R]{\textsc{Fall 2024}}
\fancyfoot[R]{
	\vspace*{\fill}\footnotesize\textcopyright{} \emph{Harvard Math 22a}	
}
\renewcommand{\headrulewidth}{0pt}

%\fancyhead[C]{}
%	\chead{}	
%	\lhead{}
%	\rhead{}
%	\cfoot{}


% Putting copyright on the footer of each page
%\fancypagestyle{secondpage}{
%	\fancyhead[C]{TESTing again} % change this to be blank
%		%\chead{}	
%	\fancyhead[L]{bears} % change this to be blank
%	\fancyhead[R]{dogs} % change this to be blank
%	\fancyfoot[R]{ %keeps the footer the same
%		\vspace*{\fill}\footnotesize\textcopyright{} \emph{Harvard Math 22a}	
%	}
%}
%\renewcommand{\headrulewidth}{0pt}
%\pagestyle{fancy}
\pagestyle{empty}


%\lhead{}
%\rhead{}
%\rfoot{\vspace*{\fill}\footnotesize\textcopyright{} \emph{Harvard Math 22a}}
%\renewcommand{\headrulewidth}{0pt}
%\pagestyle{fancy}




%%
%% Common shortcut commands:
%%
\newcommand{\ds}{\displaystyle}
\newcommand{\sgn}{\operatorname{sgn}}
\renewcommand{\Pr}{\operatorname{P}}
\newcommand{\CPr}[3][{}]{\Pr_{\text{#1}}\left(#2\,|\,#3\right)}

\newcommand{\C}{\mathbb{C}}
\newcommand{\R}{\mathbb{R}}
\newcommand{\Z}{\mathbb{Z}}
\newcommand{\N}{\mathbb{N}}
\newcommand{\Q}{\mathbb{Q}}
\newcommand{\D}{\mathbb{D}}

\newcommand{\rref}{\operatorname{rref}}
\newcommand{\im}{\operatorname{im}}
\newcommand{\rank}{\operatorname{rank}}
\newcommand{\diag}{\operatorname{diag}}
\newcommand{\Span}{\operatorname{span}}
%\renewcommand{\vec}[1]{\mathbf{#1}}
\newcommand{\charpoly}{\mathscr{P}}
\newcommand{\nicefrac}[2]{\sfrac{#1}{#2}} % using xfrac package
\newcommand{\topology}{\mathcal{T}}
\newcommand{\Inn}{\operatorname{Inn}}

\newcommand{\blankbox}[2]{\fbox{\rule{#1}{0in}\rule{0in}{#2}}}

\newenvironment{myindentpar}[1]%
 {\begin{list}{}%
         {\setlength{\leftmargin}{#1}
          \setlength{\rightmargin}{\leftmargin}}%
         \item[]%
 }
 {\end{list}}
\newtheorem{theorem}{Theorem}
\newtheorem*{NStheorem}{Theorem}
\newtheorem{cor}[theorem]{Corollary}
\newtheorem*{claim}{Claim}
\newtheorem{defn}{Definition}

\newenvironment{amatrix}[1]{%
  \left(\begin{array}{@{}*{#1}{c}|c@{}}
}{%
  \end{array}\right)
}

%--------grstep
% For denoting a Gauss' reduction step.
% Use as: \grstep{\rho_1+\rho_3} or \grstep[2\rho_5 \\ 3\rho_6]{\rho_1+\rho_3}
\newcommand{\grstep}[2][\relax]{%
   \ensuremath{\mathrel{
       {\mathop{\longrightarrow}\limits^{#2\mathstrut}_{
                                     \begin{subarray}{l} #1 \end{subarray}}}}}}
\newcommand{\swap}{\leftrightarrow}


\usetikzlibrary{arrows}
\usepackage{wrapfig}

\makeatletter
\renewcommand*\env@matrix[1][*\c@MaxMatrixCols c]{%
	\hskip -\arraycolsep
	\let\@ifnextchar\new@ifnextchar
	\array{#1}}
\makeatother

\def\env@matrix{\hskip -\arraycolsep
	\let\@ifnextchar\new@ifnextchar
	\array{*\c@MaxMatrixCols c}}

\chead{\Large\textbf{Vector Spaces, $\vec\S$4.1}}% 

\begin{document}
\thispagestyle{fancy}

Recall from last class:

\begin{itemize}
\item {\bf Theorem 3.3:} If $A$ is an $n\times n$ matrix and $E$ is an elementary matrix, then $$\operatorname{det}(EA)=\operatorname{det}(E) \operatorname{det}(A),$$ and 
$$\operatorname{det}(E)=\begin{cases}
1 & \text{if $E$ represents adding a multiple of a row to another row,}\\
-1 & \text{if $E$ represents interchanging two rows,}\\
r & \text{if $E$ represents scaling a row by a factor of $r$.}
\end{cases}$$
\item {\bf Theorem 3.6:} If $A$ and $B$ are $n$ by $n$ matrices, then $$\operatorname{det}(AB)=\operatorname{det}(A)\operatorname{det}(B).$$
\item {\bf Corollary:} $A$ is invertible if and only if $\operatorname{det}(A)\neq0$.
\item {\bf Theorem:} If $A$ is a $2 \times 2$ matrix, then the area of the parallelogram determined by the columns of $A$ is $|\det{A}|$. If $A$ is a $3 \times 3$ matrix, then the volume of the parallelepiped is $|\det{A}|$.
\item {\bf Theorem:} Let $T: \mathbb{R}^2 \to \mathbb{R}^2$ be linear with associated matrix $A$. If $S$ is a parallelogram in $\mathbb{R}^2$, then $\text{Area}(T(S))=|\det{A}|\text{Area}(S)$. 
\end{itemize}

\newpage

\begin{defn}
A real {\bf vector space} is a non-empty set $V$ on which two operations, scalar multiplication and addition, are defined subject to the following 10 axioms. We call elements in $V$ vectors. Let $u,v,w \in V$ and $c,d\in \mathbb{R}$.
\begin{enumerate}
\item $u+v \in V$.
\item $u+v=v+u$.
\item $(u+v)+w=u+(v+w)$.
\item there exists a vector $0_{V}$ in $V$ such that $u+0_{V}=u$.
\item there exists a vector $(-u)$ in $V$ such that $u+(-u)=0_{V}$.
\item $cu \in V$.
\item $c(u+v)=cu+cv$.
\item $(c+d)u=cu+du$.
\item $c(du)=(cd)u$.
\item $1\cdot u =u$.
\end{enumerate}
\end{defn}

\begin{enumerate}

\Problem Determine if the following are vector spaces.

\begin{enumerate}
\item $\mathbb{R}^n$ with the usual operations.

\vfill
\item $S=\{(y_0, y_1, \dots ): y_k \in \mathbb{R} \text{ for  all } k \geq 0\}$. Define addition component-wise and scalar multiplication by scaling every component.

\vfill

\newpage

\item Let $V= \left\{ \begin{bmatrix} x\\ y\\ \end{bmatrix} \in \mathbb{R}^2 : x^2+y^2 =1 \right \}$ with the usual operations.

\vfill
\item Let $\mathbb{P}_n =\{ a_0+a_1 t +\dots + a_n t^n : a_0, \dots, a_n \in \mathbb{R} \}$ with operations defined by corresponding operations on polynomials.

\vfill
\item Let $M_{n\times n}$ be the set of $n\times n$ matrices with the operations coming from matrix addition and scaling a matrix.

\vfill

\item Let $W=\{(x,y,1): x,y\in \mathbb{R} \}$ with the usual operations.

\vfill

\item Let $Z=\{(a,b) : a,b \in \mathbb{R} \}$ and $(a,b)+(c,d)=(a+c,0)$ and $r(a,b)=(ra,0)$.

\vfill

\item Let $Y=\{(a,b) : a,b \in \mathbb{R} \}$ and $(a,b)+(c,d)=(a+c,b-d)$ and $r(a,b)=(ra,rb)$.

\vfill

\item Let $F=\{ f: \mathbb{R} \to \mathbb{R} : f \text{ is a function} \}$ with operations coming from the usual operations on functions.

\vfill
\newpage

\end{enumerate}

\begin{defn}
A {\bf subsapce} of a vector space $V$ is a subset of $V$ such that
\begin{enumerate}
\item $0_V \in H$.
\item $H$ is closed under vector addition; i.e. if $u, v \in H$, then $u+v \in H$.
\item $H$ is closed under scalar multiplication; i.e. if $c \in \mathbb{R}$ and $u \in H$, then $cu \in H$.
\end{enumerate}
\end{defn}

\Problem Let $V$ be a vector space.
\begin{enumerate}
\item Is $H=\{ 0_V \}$ a subspace of $V$?

\vfill

\item Suppose $V=\mathbb{R}^3$, then is $\mathbb{R}^2$ a subspace of $V$?

\vfill

\end{enumerate}

{\bf Theorem 4.1} If $v_1, \dots, v_p \in V$, then $\text{Span}\{v_1 ,\dots, v_p\}$ is a subspace of $V$. We call it the subspace generated by or spanned by $v_1,\dots, v_p$.

\Problem Prove the Theorem.

\vfill

\newpage

\Problem Consider the vector space of polynomials of degree 2 or less. We denote this vector space by $\mathbb{P}_2$.

\begin{enumerate}
\item Is the set of polynomials of degree 2 a subspace?

\vfill

\item Let $u=1+t$, $v=t$, and $w=1+t+2t^2$. Is $w \in \text{Span}\{ u, v\}$?

\vfill

\end{enumerate}

\Problem Let $W_1$ and $W_2$ be subspaces of a vector space $V$. 

\begin{enumerate}
\item When is $W_1 \cap W_2$ a subspace of $V$?

\vfill

\item When is $W_1 \cup W_2$ a subspace of $V$?

\vfill



\end{enumerate}

\Problem Jonathan's question: If $v$ is a vector space, then is $(-1) u= -u$?

\vfill

\end{enumerate}

\begin{comment}
\newpage
\chead{\Large\textbf{Vector Spaces -- Answers / Solutions}}%
\lhead{}
\rhead{}
\thispagestyle{fancy}

\begin{enumerate}
  \Answer 
  
  \begin{enumerate}
  	\item We have previous seen that $\mathbb{R}^n$ with vector addition and scalar multiplication satisfies the 10 axioms, and hence it is a vector space.
  	\item $S$ is a vector space.
  	\item $V$ is not a vector space, since adding two elements of $V$ doesn't always yield another element in $V$.
  	\item $\mathbb{P}_n$ is a vector space.
  	\item $M_{n\times n}$ is a vector space.
  	\item $W$ is not closed under scalar multiplication or vector addition, so it is not a vector space.
  	\item $Z$ does not have a zero vector, so it is not a vector space.
  	\item The addition operation in $Y$ is not commutative, so $Y$ is not a vector space.
  	\item $F$ is a vector space.
  \end{enumerate}

\Answer
\begin{enumerate}
	\item Yes, $H$ is a subspace. We can prove it be checking the 3 conditions to be a subspace.
	
	{\bf (a)} Since $0_V \in H$, we know (a) holds.
	
	{\bf (b)} By vector space axiom 4, we know $0_V+0_V = 0_V$. Thus $H$ is closed under vector addition.
	
	{\bf (c)} Let $c \in \mathbb{R}$, then we want to prove $c 0_V \in H$. By axiom 7, $c 0_V = c (0_V +0_v)= c 0_V +c 0_V$. Now subtracting $c 0_V$ from both sides (axiom 5) yields $c 0_V=0_V$, and hence $H$ is closed under scalar multiplication.
	
	\item No, $\mathbb{R}^2$ is not a subset of $V$, so it cannot be a subspace of $V$.
\end{enumerate}
\Answer

\begin{proof}
	We need to check the three subspace conditions.
	
	{\bf (a)} Since $0_V= 0\cdot v_1+\dots +0 \cdot v_p$, we know $0_V \in \text{Span}\{v_1,\dots, v_p\}$.
	
	{\bf (b)} Let $u,v \in \text{Span}\{v_1, \dots, v_p\}$. By definition of the span, we know there exists scalars $c_1,\dots c_p, d_1, \dots, d_p\in \mathbb{R}$ such that $$u=c_1 v_1+\dots +c_p v_p$$ and $$v=d_1 v_1+\dots + d_p v_p.$$ By vector space axioms (2) and (8), we get 
	\begin{align*}
		u+v &= c_1 v_1+\dots +c_p v_p+ d_1 v_1+\dots + d_p v_p\\
		&= c_1 v_1 +d_1 v_1 +\dots + c_p v_p +d_p v_p\\
		&= (c_1+d_1)v_1+\dots + (c_p+d_p)v_p.
	\end{align*}
Thus $u+v$ is a linear combination of $v_1, \dots, v_p$ and $u+v \in \text{Span}\{v_1,\dots, v_p\}$.

{\bf (c)} Let $c \in \mathbb{R}$ and $u \in \text{Span}\{v_1,\dots, v_p\}$. Again by definition of span, there exists $c_1, \dots, c_p \in \mathbb{R}$ such that $u=c_1 v_1 +\dots + c_p v_p$, and using axioms 7 and 9, we get $$cu= c (c_1 v_1+\dots+ c_pv_p)=(c c_1) v_1 +\dots + (c c_p) v_p.$$ Therefore $cu \in \text{Span}\{v_1, \dots, v_p\}$.
\end{proof}

\Answer
\begin{enumerate}
	\item No. The zero vector in $\mathbb{P}_2$ is the zero polynomial; namely, $0_{\mathbb{P}_n}=0 + 0 \cdot t + 0 \cdot t^2$. However, the zero polynomial is not of degree 2, so the set of polynomials of degree 2 fails the first condition to be a subspace.
	\item No. For $w$ to be in the span of $u$ and $v$, we would need to find scalars $c_1, c_2$ such that $w=c_1 u+c_2 v$, or equivalently $1+t+2t^2=c_1(1+t)+c_2t$. This is impossible since the left hand side is quadratic and the right hand side is linear. Thus $w\not \in \text{Span}\{u,v\}$.
\end{enumerate}
\Answer
\begin{enumerate}
	\item $W_1 \cap W_2$ is a subspace of $V$.
	
	\begin{proof}
	{\bf (a)} Since $W_1$ and $W_2$ are both subspaces of $V$, we know $0_V \in W_1$ and $0_V \in W_2$. Thus $0_V \in W_1 \cap W_2$.
	
	{\bf (b)} Let $u, v \in W_1\cap W_2$. Since $u, v \in W_1$ and $W_1$ is a subspace, we know $u+v \in W_1$. Since $u, v \in W_2$ and $W_2$ is a subspace, we know $u+v \in W_2$. Thus $u+v \in W_1 \cap W_2$.
	
	{\bf (c)} Let $c \in \mathbb{R}$ and $u \in W_1 \cap W_2$. Since $W_1$ is a subspace, we get $cu \in W_1$. Since $W_2$ is a subspace, we get $cu \in W_2$. Thus $cu \in W_1 \cap W_2$.
	
	Thus $W_1 \cap W_2$ is a subspaace of $V$.
	\end{proof}
	
	\item First note that generally $W_1 \cup W_2$ is not subspace of $V$. For instance take $V=\mathbb{R}^2$ and $W_1=\text{Span}\left \{\begin{bmatrix} 1\\ 0\\ \end{bmatrix} \right\}$ and $W_2=\text{Span}\left\{\begin{bmatrix} 0\\ 1\\ \end{bmatrix} \right\}$. Then $\begin{bmatrix} 1\\ 0\\ \end{bmatrix}+ \begin{bmatrix} 0\\ 1\\ \end{bmatrix} \not \in W_1 \cup W_2$.
	
	{\bf Claim:} $W_1 \cup W_2$ is a subspace of $V$ if and only if $W_1 \subseteq W_2$ or $W_2 \subseteq W_1$.
	
	\begin{proof} We prove each implication directly.
		
	$(\implies)$ Let $x\in W_1$ and  $y \in W_2$. Since $W_1 \cup W_2$ is a subspace, we know $x+y \in W_1 \cup W_2$.
	
	{\bf Case 1:} $x+y \in W_1$.
	
	Now $y=(x+y)-x$, so $y$ is the sum of two elements in $W_1$. Since $W_1$ is a subspace, we get $y \in W_1$. Thus $W_2 \subseteq W_1$.
	
	{\bf Case 2:} $x+y \in W_2$.
	
	Here $x=(x+y)-y$, so $x$ is the sum of two elements in $W_2$. Since $W_2$ is a subspace, we get $x \in W_2$. Thus $W_1 \subseteq W_2$.
	
	Therefore, if $W_1 \cup W_2$ is a subspace of $V$, then $W_1 \subseteq W_2$ or $W_2 \subseteq W_1$.
	
	$(\impliedby)$ If $W_1 \subseteq W_2$, then $W_1 \cup W_2=W_2$, and hence $W_1 \cup W_2$ is a subspace of $V$. If $w_2 \subseteq W_1$, then $W_1 \cup W_2 = W_1$, and hence $W_2 \cup W_2$ is a subspace of $V$. In either case, $W_1 \cup W_2$ is a subspace of $V$.
	\end{proof}
\end{enumerate}
\item {\bf Jonathan's Claim: } If $V$ is a vector space, then $(-1) u = -u$.

\begin{proof}
Note that by definition $-u$ is the additive inverse of $u$.

{\bf Subclaim:} additive inverses are unique.

{\bf Proof of subclaim:} Suppose $x$ and $y$ are both additive inverses for $u$; namely $u+x=0$ and $u+y=0$. Now we have $$x= 0+x = (u+y)+x = y+(u+x) = y+0 =y$$ where we used axioms 2, 4, 3. 
 
Note $$u+ (-1)u = (1 +(-1)) u = 0 \cdot u =0.$$ Thus $(-1)u$ is the additive inverse of $u$; namely, $(-1)u =-u$.

\end{proof}
\end{enumerate}  

\end{comment}


\end{document}


%%% Local ables: 
%%% mode: latex
%%% TeX-master: t
%%% End: 
